\begin{abcliste}
\abc
\result{\text{Paar 2 ist falsch.}} Beide Kräfte wirken auf denselben Körper, die Lampe. Kraft und Gegenkraft wirken aber immer auf zwei verschiedene Körper: Sie gehören zu einer Wechselwirkung, „A auf B“ $\leftrightarrow$ „B auf A“. Dass sich die beiden Kräfte an der Lampe aufheben, ist das erste Newtonsche Gesetz (die Lampe ruht), nicht das dritte. In den Paaren 1, 3 und 4 sind die beiden Körper vertauscht; sie stimmen.

\begin{center}
\begin{tikzpicture}[>=Stealth]
 \fill[pattern=north east lines] (-1.5,3) rectangle (1.5,3.25);
 \draw[thick] (-1.5,3) -- (1.5,3);
 \draw[thick] (0,3) -- (0,1.2);
 \draw[fill=gray!30] (-0.25,1.2) -- (0.25,1.2) -- (0.55,0.5) -- (-0.55,0.5) -- cycle;
 \draw[->, very thick] (-0.12,0.85) -- (-0.12,-0.3) node[left] {Erde $\to$ Lampe};
 \draw[->, very thick] (-0.12,1.2) -- (-0.12,2.4) node[left] {Kabel $\to$ Lampe};
 \draw[->, very thick, gray] (0.12,1.2) -- (0.12,0.0) node[right] {Lampe $\to$ Kabel};
 \draw[thick] (-3,-1.2) -- (3.6,-1.2);
 \fill[pattern=north east lines] (-3,-1.45) rectangle (3.6,-1.2);
 \draw[->, very thick, gray] (3.0,-1.2) -- (3.0,-0.4) node[right] {Lampe $\to$ Erde};
\end{tikzpicture}
\end{center}
(Schwarz: die beiden Kräfte von Paar 2, beide auf die Lampe; grau: ihre richtigen Gegenkräfte. Der graue Pfeil rechts unten wirkt auf die Erde.)

\abc
Die Gegenkraft zu „Die Erde zieht die Lampe nach unten“ ist \result{\text{„Die Lampe zieht die Erde nach oben“}}. Die Gegenkraft zu „Das Kabel zieht die Lampe nach oben“ ist \result{\text{„Die Lampe zieht das Kabel nach unten“}}. Beide Paare stehen schon in Bens Liste (1 und 4); Paar 2 verbindet zwei Kräfte aus zwei verschiedenen Wechselwirkungen.

\abc
Die Lampe ruht, also hält das Kabel der Gewichtskraft das Gleichgewicht:
\begin{align}
 F_\mathrm{G} &= \FGF = \mL\cdot\g = \FG \approx \result{\FGP{0}{3}}
\end{align}
Gegenkräfte sind gleich gross, und das leichte Kabel ruht ebenfalls (die Decke zieht es so stark nach oben, wie die Lampe es nach unten zieht). Darum sind \result{\text{alle Kräfte in Bens Liste } \FGP{0}{3}} gross. Dass die Kräfte von Paar 2 gleich gross sind, macht sie trotzdem nicht zu einem Paar.
\end{abcliste}
